\section{Part A. Virtual surface structures and realization}

We use $S_{g,n}$ for a connected orientable surface obtained by deleting $n$
points from a compact surface of genus $g$, and put
$P_{g,n}=\pi_1(S_{g,n},*)$. In this part ``hyperbolic'' means
$2g-2+n>0$. A peripheral structure consists of the conjugacy classes of the
puncture subgroups, their primitive parametrizations by a common orientation
line $T\cong\mathbb Z$, and its identification with the surface orientation
line. No basis of $T$ is chosen in the full structure: an automorphism
may act on it by $-1$. Choosing a basis defines the superscript $+$.
Homeomorphisms may permute punctures. Our $E,P,Q$
correspond respectively to Grothendieck's $\Pi,\Pi_0,G$ in \S52.1,
folios 737--739, scans 179--183. His later $T$ and $k_i^\natural$ occur in
folios 741--743, scans 188--191.

\subsection{Definitions, and a necessary correction to the LT definition}

\begin{statement}{A.1}{Proved}{A precise discrete L definition and the meaning of the headings}
For a discrete group $E$, let $\mathcal L_0(P)$ be the set of hyperbolic
surface peripheral and orientation structures on $P$. Then the definition
in \S52.1, folios 737--738, scans 180--181, is
\[
 \mathcal L(E)=\mathop{\mathrm{colim}}_{P\triangleleft_f E}
             \mathcal L_0(P)^{E/P}.
\]
The transition maps are restriction to finite coverings; an L-group is
$(E,\lambda)$ with $\lambda\in\mathcal L(E)$. Each transition map is
injective. The alternative description by germs in scans 188--189 uses
outer homomorphisms $nT\to E$, two representatives being identified when
their restrictions to a common finite-index subgroup of $T$ are conjugate.
The table's $\Pi L$ heading in scan 4 supplies no third, independent
definition: the body calls the proposed enlargement $LT$ and also $TL$.
\end{statement}
\begin{proof}
An oriented puncture subgroup in a hyperbolic surface group is maximal
cyclic; its positive finite-index subgroups determine it and its positive
generator uniquely. Indeed, a nontrivial element acts on the universal
hyperbolic plane with an axis or a unique parabolic fixed point; its
centralizer in the surface group is the corresponding cyclic subgroup.
The same argument shows unique roots. If two peripheral structures have
the same pullback to a finite-index subgroup, taking the maximal cyclic
overgroups recovers the same peripheral classes upstairs. Orientation is
recovered by the positive degree of a finite covering (equivalently by
the relative fundamental class after replacing punctures by boundary
circles). Thus restriction is injective. A common normal core makes the
indexing system filtered. This proves that the displayed colimit is an
honest set of virtual structures, with no additional quotient ambiguity.
The germ construction records precisely the finite-index intersections
of those peripheral subgroups. For a closed defining surface, the
orientation class must still be included; an empty list of homomorphisms
alone does not encode it.

The transcription of scan 180 does not legibly retain all four excluded
types. The mathematical assertion about ``anab\'elien'' requires exactly
$(0,0),(0,1),(0,2),(1,0)$ to be excluded. This is evidence for an editorial
correction, not a claim to have deciphered the missing scan. We inspected
the supplied scan 4: its earlier heading is struck, and the body at
scans 190--191 explicitly explains $L$ as ``lacets'' and $T$ as
``Teichm\"uller''. We therefore preserve $\Pi L$ in the reading ledger
without inventing axioms for another object.
\end{proof}

\begin{statement}{A.2}{Refuted}{Germs alone do not determine the LT defining subgroup}
The assertion in \S52.1, folios 742--743, scans 190--191, that the
normal defining subgroup is determined up to commensurability does not
follow from the stated germ-only axiom. A corrected LT datum includes
the defining commensurability class as part of its data. The same
counterexample applies in the profinite category.
\end{statement}
\begin{proof}
Take $E=P_{2,0}\times P_{2,0}$, give $T=\mathbb Z$ its chosen orientation,
and take the empty peripheral list. Both
$P_{2,0}\times1$ and $1\times P_{2,0}$ are normal subgroups on which the
induced empty peripheral list is a closed hyperbolic surface structure.
Their intersection is trivial and has infinite index in each. Thus both
satisfy the written existence axiom but they are not commensurable.
Replacing every group by its profinite completion gives the same
example: the trivial intersection is not open in either factor.

We repair the definition by specifying a commensurability class
$\mathscr P$ of normal defining subgroups, compatible oriented surface
structures on its members, and the peripheral germ data. In the
discrete category commensurable means that the intersection has finite
index in both groups. In the profinite category the groups are closed,
commensurable means that the intersection is open in both, all maps are
continuous, and $T$ is a free rank-one $\widehat{\mathbb Z}$-module.
The existence of one normal defining subgroup is retained. This is
also the structure actually available in an arithmetic extension:
the geometric kernel is specified, rather than reconstructed from an
empty list of punctures. With this datum, any extension by a group
preserving the specified virtual surface structure is an LT-group,
as in scan 191.
\end{proof}

\subsection{Realization and the finite ineffective kernel}

\begin{statement}{A.3}{Known}{Every discrete hyperbolic L-group is realized by a finite surface action}
Let $(E,\lambda)$ be a discrete L-group with a normal defining surface
group $P\cong P_{g,n}$ and put $Q=E/P$. There is an action of the finite
group $Q$ on $S_{g,n}$ by diffeomorphisms, possibly orientation reversing
and possibly ineffective, and an isomorphism of extensions and
L-structures
\[
 E\ \cong\ \pi_1(S_{g,n},Q;*).
\]
Here the right side consists of pairs $(q,[\alpha])$, with $\alpha$ a
path from $*$ to $q(*)$, and product induced by concatenation and the
$Q$-action. This settles the existence question in \S52.1, folio 739,
scan 183, in the hyperbolic meaning of A.1.
\end{statement}
\begin{proof}
Conjugation gives a homomorphism
$\rho:Q\to\operatorname{Out}_L(P)$, independent of representatives
in $E$. The Dehn--Nielsen--Baer theorem used here identifies
$\operatorname{Out}_L(P)$ with isotopy classes of all
homeomorphisms of $S_{g,n}$, permitting puncture permutations:
an outer automorphism preserving the peripheral classes with
one common orientation sign is induced by a homeomorphism,
unique up to isotopy. For surfaces with punctures, replace
punctures by boundary circles and apply the relative-boundary
form \citep[Theorem 2.3]{BeringQingWigglesworth2024};
for closed surfaces the identification is also explicitly
stated on p.~235 of \citet{Kerckhoff1983}.
This sends $\rho$ to a finite
subgroup $H=\rho(Q)$ of the extended mapping class group of $S_{g,n}$.
The relevant theorem is finite subgroup realization:
\citet[Theorem 5 and its finite-type extension in \S4; Theorem 7]{Kerckhoff1983}
realizes $H$ by isometries of a complete finite-area hyperbolic metric.
That paper explicitly permits orientation-reversing classes. Let $Q$
act through $Q\to H$.

For completeness, the extension comparison here is algebraic and
requires no chosen splitting. Because $Z(P)=1$, the pullback
\[
 E_\rho=\{(a,q)\in\operatorname{Aut}(P)\times Q:
                  [a]=\rho(q)\}
\]
has kernel $\operatorname{Inn}(P)\cong P$. The map
$e\mapsto(\operatorname{conj}_e,\overline e)$ from $E$ to $E_\rho$
is injective: its kernel lies in $P$ and centralizes $P$. It is
surjective: for $(a,q)$ choose a lift $e$ of $q$; then
$a\operatorname{conj}_e^{-1}=\operatorname{conj}_p$ for a unique
$p\in P$, and $pe$ maps to $(a,q)$. The orbifold path extension
of the realized $Q$-action induces the same $\rho$, so the same
construction identifies it with $E_\rho$. The resulting isomorphism
restricts to the chosen identification of $P$ and therefore preserves
the virtual peripheral and orientation structure. It is uniquely
determined once that identification and the quotient $Q$ are fixed.
This proves the required reduction to the cited theorem.
\end{proof}

\begin{statement}{A.4}{Proved}{The exact centralizer criterion}
In A.3 let $K=C_E(P)$. Projection to $Q$ identifies $K$ with
$\ker\rho$; hence $K$ is finite and is the ineffective kernel of the
realized action. The following conditions are equivalent:
\[
 K=1;\qquad \rho\text{ is injective};\qquad
 C_E(J)=1\quad\text{for every finite-index subgroup }J\leq E.
\]
Consequently the faithful version anticipated in \S52.1, folio 739,
scans 183--184, is exactly correct in hyperbolic type. Triviality of
$Z(E)$ alone is insufficient.
\end{statement}
\begin{proof}
The pullback model in A.3 identifies $C_E(P)$ with the pairs $(1,q)$
for $q\in\ker\rho$. If $P'\leq P$ has finite index and
$e\in E$ centralizes $P'$, then for each $x\in P$ some positive
power $x^d$ lies in $P'$. Thus
$(exe^{-1})^d=x^d$, and unique roots in $P$ give $exe^{-1}=x$.
It follows that $C_E(P')=K$. If $J$ has finite index in $E$, then
$J\cap P$ has finite index in $P$ and $C_E(J)\subseteq K$.
Thus $K=1$ implies the centralizer condition; the reverse implication
uses $J=P$. Finally $E=P_{2,0}\times S_3$ has trivial center but
$C_E(P_{2,0})=S_3$. It is realized by the trivial $S_3$-action on
$S_{2,0}$ and is not faithful.
\end{proof}

\begin{statement}{A.5}{Proved}{Elementary exceptional types}
The centralizer criterion and the center-free uniqueness argument are
not statements about the excluded types. In particular a faithful
finite translation action on a torus has a lift group centralizing
its lattice. Nevertheless every extension of a lattice
$\Lambda\cong\mathbb Z^d$, $d=1,2$, by a finite group with the
prescribed linear action is realized by a finite affine action on
$\mathbb R^d/\Lambda$. For $d=1$ this yields the corresponding
annular action after taking a product with an interval.
\end{statement}
\begin{proof}
Choose a section of the quotient and write the extension law using
an integral two-cocycle $c(q,r)\in\Lambda$ and a representation
$A:Q\to\operatorname{GL}(\Lambda)$. Over $\mathbb R$, the
positive-degree cohomology of a finite group vanishes by averaging,
so choose $b_q\in\Lambda\otimes\mathbb R$ with
$b_q+A_qb_r-b_{qr}=c(q,r)$. Then
$q\cdot[z]=[A_qz+b_q]$ is a finite group action and its labelled
lift extension has precisely the displayed cocycle. The averaging
assertion can be checked directly on the usual inhomogeneous
cochain complex using $|Q|^{-1}\sum_{q\in Q}q$.
For the torus, translations commute with lattice translations, so
the claimed faithful-action centralizer equivalence fails there.
For $(0,0)$ and $(0,1)$ the fundamental group is trivial and cannot
recover a puncture from a nontrivial cyclic subgroup. This also
corrects the isolated phrase in scan 180 associating $\mathbb Z$
to $(0,1)$: $\mathbb Z$ belongs to $(0,2)$.
\end{proof}

\subsection{Non-uniqueness and Teichm\"uller groups}

\begin{statement}{A.6}{Proved}{Precisely which sets are torsors}
Fix a realized discrete L-group $E_X$ isomorphic to $E$. The set
$\operatorname{Isom}_L(E,E_X)$ is a right
$\operatorname{Aut}_L(E)$-torsor; after quotienting target inner
automorphisms it is a right $\operatorname{Out}_L(E)$-torsor.
The set of conformal realizations is instead a union of orbit spaces.
For a fixed effective orbifold signature $\sigma$ it is
$\mathcal T_\sigma/\operatorname{Mod}_\sigma$; allowing all defining
finite-index surface subgroups gives the union over their corresponding
covering data. Thus an unspecified ``set of realizations'' is not
canonically a single torsor.
\end{statement}
\begin{proof}
If $u,v:E\to E_X$ are isomorphisms, the unique element sending
$u$ to $v$ by precomposition is $u^{-1}v$. Modulo target conjugation,
this identifies the ambiguity with outer automorphisms. If orientation
is fixed, replace each group by its $+$ subgroup.

For the geometric statement, a marking of an orbifold is a
homeomorphism from a fixed topological orbifold, considered up to
isotopy; removing the marking is exactly the mapping class action.
Each normal torsion-free subgroup of its orbifold group gives a
finite surface cover; pulling back a conformal structure gives the
corresponding finite action, and descending the invariant structure
is inverse to pullback. This proves the orbit-space statement with
all objects and equivalences specified.

Three explicit examples make the ambiguity concrete. For the
orbifold $(0;2,3;1)$, whose group is $C_2*C_3$, the oriented
Teichm\"uller space is a point and its oriented outer L-group is
trivial. The same is true of $(0;2,4;1)$ with group $C_2*C_4$.
For $S_{1,1}$ the oriented marking ambiguity is
$\operatorname{SL}_2(\mathbb Z)$ and the conformal realization space
is $\mathbb H/\operatorname{SL}_2(\mathbb Z)$. Indeed a pointed
complex torus is $\mathbb C/(\mathbb Z+\tau\mathbb Z)$,
$\operatorname{Im}\tau>0$, and changing an oriented basis acts
by fractional linear transformations. With the order-four symmetry
specified, the equation $\tau=-1/\tau$ leaves the unique point
$\tau=i$. In the first two examples the full outer L-group has
order two, represented by reflection. These are computations of
both marking and geometric ambiguity, not assertions that a
positive-dimensional moduli space is a group torsor.
\end{proof}

\begin{statement}{A.7}{Refuted}{The proposed characterization of the modular group}
The conjunction of trivial oriented outer L-group and trivial
centralizers of all finite-index subgroups does not characterize
$\operatorname{PSL}_2(\mathbb Z)$, contrary to the tentative
question in \S52.1, folios 740--741, scans 185--188.
The L-group $C_2*C_4$ of signature $(0;2,4;1)$ is a counterexample.
\end{statement}
\begin{proof}
Write $E=\langle a,b\mid a^2=b^4=1\rangle$ and take the cusp
class represented by $(ab)^{-1}$. The map $E\to C_4$ with
$a\mapsto2$, $b\mapsto1$ is surjective and is injective on both
finite factors. The kernel is torsion-free because every torsion
element of a free product of finite groups is conjugate into a
factor (a cyclically reduced word of length at least two has
infinite order). The index-four cover has Euler characteristic
$4(\tfrac12+\tfrac14-1)=-1$. The cusp image has order four, so
there is one cusp upstairs; hence the surface is $S_{1,1}$.
This proves the L assertion, including the actual finite surface
action of order four.

The orbifold has one cusp and two cone points of distinct orders.
Every oriented orbifold mapping class fixes all three distinguished
points and is isotopic to the identity: cut the sphere along two
arcs joining them and use the Alexander isotopy on the resulting
disc. The orbifold Dehn--Nielsen--Baer theorem identifies this
group with $\operatorname{Out}_L^+(E)$; the peripheral and rotation
classes are the allowable classes of
\citet[\S4, pp.~255--256]{Kerckhoff1983}. Thus it is trivial.
An equivalent algebraic verification uses the standard free-product
normal form: an automorphism of $C_2*C_4$ is, modulo inner
automorphism, induced by automorphisms of the two factors. The
only nonidentity option sends $b$ to $b^{-1}$ and reverses the
cusp orientation. Its centralizers of finite-index subgroups
are trivial by A.4. Finally
$E^{\mathrm{ab}}=C_2\times C_4$ has order eight, whereas
$(C_2*C_3)^{\mathrm{ab}}=C_2\times C_3$ has order six, so the
groups are not isomorphic. The finite quotient, element orders,
Euler characteristic and abelianization distinction are checked
by \texttt{code/a\_small\_actions.py}.

The valid replacement is a classification by the complete orbifold
signature and peripheral orientation data. Rigidity of the
oriented three-pointed sphere explains why many signatures, not
just $(0;2,3;1)$, have trivial oriented outer group.
\end{proof}

\begin{statement}{A.8}{Proved}{The modular-group statements themselves}
For $E=C_2*C_3$ with its cusp structure,
$\operatorname{Out}_L^+(E)=1$,
$\operatorname{Out}_L(E)=C_2$, and conjugation gives
$\operatorname{PGL}_2(\mathbb Z)\cong\operatorname{Aut}_L(E)$.
Every surjection $E\to S_3$ is conjugate by an element of $S_3$
to the standard one; consequently its kernel $F_2$ is
characteristic. These are the assertions, as distinct from the
failed characterization, of scans 185--188.
\end{statement}
\begin{proof}
The oriented mapping class computation is exactly the first
three-pointed orbifold computation in A.6. Reflection supplies
the nontrivial outer class, and the absence of a centralizer in
A.4 proves injectivity of the conjugation map from
$\operatorname{PGL}_2(\mathbb Z)$. Its image contains all inner
automorphisms and reflection, so it is surjective.
For the last assertion let $r^3=s^2=1$ generate $E$.
Surjectivity forces the image of $r$ to be one of the two
3-cycles and that of $s$ to be one of the three transpositions.
There are six pairs. Conjugation by $S_3$ acts freely on them,
since a commuting element centralizes their generated group,
whose center is trivial; hence there is one orbit. Precomposing
the standard surjection with any automorphism gives another
surjection with the same kernel. The code enumerates all six
pairs and their one conjugation orbit.
\end{proof}

\begin{statement}{A.9}{Proved}{Ineffective L-groups require stack mapping classes}
For a hyperbolic surface group $P$ and a finite abelian group $A$,
give $P\times A$ its induced virtual L-structure. Then
\[
 \operatorname{Out}_L(P\times A)
 \cong\operatorname{Hom}(P^{\mathrm{ab}},A)
 \rtimes\bigl(\operatorname{Out}_L(P)\times\operatorname{Aut}(A)\bigr).
\]
Thus identifying every L-group's Teichm\"uller group with the
mapping class group of its effective coarse orbifold would be
false. The correct object is the mapping class group of the
possibly ineffective quotient stack, retaining its gerbe data.
\end{statement}
\begin{proof}
The maximal finite normal subgroup of $P\times A$ is $A$:
projection of any finite subgroup to the torsion-free group $P$
is trivial. It is therefore characteristic. Every automorphism
has the unique form
$(p,a)\mapsto(\phi(p),\chi(p)+\eta(a))$ with
$\phi\in\operatorname{Aut}(P)$,
$\eta\in\operatorname{Aut}(A)$ and
$\chi\in\operatorname{Hom}(P,A)$. It preserves the virtual
L-structure exactly when $\phi$ preserves that of $P$;
the finite component of a peripheral image disappears upon
restriction to a multiple of $|A|T$. Quotienting inner
automorphisms replaces $\operatorname{Aut}_L(P)$ by
$\operatorname{Out}_L(P)$ and leaves $\chi$ unchanged.
The displayed semidirect action is precomposition by $\phi^{-1}$
and postcomposition by $\eta$. For example $P=P_{2,0}$ and
$A=C_2$ contribute an additional subgroup $(C_2)^4$ acting
trivially on the coarse surface. A coarse-orbifold mapping class
group cannot record this kernel.
\end{proof}

\subsection{Finite computations and the profinite boundary of the argument}

\begin{statement}{A.10}{Proved}{What the finite enumeration certifies}
The attached program constructs multiplication tables for all 24
abstract groups of order at most 12 and checks their group axioms.
It checks every Riemann--Hurwitz-admissible signature for a
faithful orientation-preserving action on a closed genus-two
surface: 42 group/signature rows, of which 17 admit a generating
vector. It also gives 126 rows obtained from nontrivial normal
ineffective kernels in these groups. In the separately specified
box $0\leq h\leq2$, $0\leq s\leq3$, at most three cone points
of orders at most 12, it lists 4,368 signatures, of which 4,258
are hyperbolic and 66 are bad orbifolds.
\end{statement}
\begin{proof}
For a group $Q$ of order $d$, an oriented action on a closed
surface of genus two is specified by an epimorphism from
\[
 \bigl\langle a_i,b_i,c_j\mid c_j^{m_j}=1,\quad
       \prod_{i=1}^{h}[a_i,b_i]\prod_jc_j=1\bigr\rangle
       \longrightarrow Q
\]
that sends each $c_j$ to an element of exact order $m_j$, with
$2/d=2h-2+\sum_j(1-1/m_j)$. Exact orders make its kernel
torsion-free. Conversely the branched-cover monodromy of every
such action gives this datum. For $d\leq12$ the equation forces
$h\leq2$ and at most six cone points, and each $m_j$ divides $d$.
The program enumerates precisely these possibilities; it tests
all handle images and all but the last cone image, which is
forced by the relation. Closure under multiplication checks
generation. This is an exhaustive existence test, not a count
of Hurwitz orbits. All normal subgroups of each finite group
are generated by closure, so pullback from their quotients
accounts for ineffective actions with these bounds. A.3 proves
that no separate nongeometric extension with the same outer
action is missing.

The noncompact signature table is a finite box with the bounds
explicitly stated above. The phrase ``genus at most two'' alone
does not impose a bound on the number of punctures: even
$P_{0,n}$ with the trivial acting group gives infinitely many
L-groups. No finite table can be an exhaustive listing under
that unbounded interpretation. Closed genus zero and one do
not meet the hyperbolic L axiom; their elementary realization
was treated in A.5. Orientation-reversing actions are covered
by A.3 but are not included in this oriented signature count.
\end{proof}

\begin{statement}{A.11}{Proved}{An exact profinite realization criterion with discretization}
Let $1\to\widehat P\to E\to Q\to1$ be an extension of
profinite groups, with $Q$ finite and $\widehat P$ the full
profinite completion of a specified hyperbolic surface group
$P$, including its completed loop structure. Suppose the outer
action $Q\to\operatorname{Out}_L(\widehat P)$ factors through
$\operatorname{Out}_L(P)$. Then $E$ is the full profinite
completion of the discrete orbifold path extension constructed
in A.3. Conversely a realization preserving this specified
discretization has that factorization. All isomorphisms here
are continuous and the comparison is an isomorphism of
extensions with the specified kernel and quotient.
\end{statement}
\begin{proof}
Construct the discrete extension $E_0$ from the lifted outer
action by A.3. The profinite topology induced on its
finite-index subgroup $P$ is the full one: for any finite-index
subgroup $N\leq P$, its core in $E_0$ is finite index and is
contained in $N$. Hence completion gives an exact sequence
$1\to\widehat P\to\widehat E_0\to Q\to1$.
The geometric surface completion $\widehat P$ is center-free;
here is a proof that does not treat the earlier LM3 return as
authority. If $z$ is a nonidentity central element, choose an
open normal subgroup $\widehat N$ whose quotient $D$ detects
$z$, and put $N=P\cap\widehat N$. The cellular chain complex
of the regular $D$-cover gives, over $\mathbb Q$,
\[
 H_1(N,\mathbb Q)\cong
 \begin{cases}
  \mathbb Q\oplus\mathbb Q[D]^{r-1},& P=F_r,\ r\geq2,\\
  \mathbb Q^2\oplus\mathbb Q[D]^{2g-2},& P=\pi_1(S_{g,0}),\ g\geq2.
 \end{cases}
\]
Indeed, in the graph case $C_1=\mathbb Q[D]^r$,
$C_0=\mathbb Q[D]$, $H_0=\mathbb Q$; in the closed case
$C_0=C_2=\mathbb Q[D]$, $C_1=\mathbb Q[D]^{2g}$ and
$H_0=H_2=\mathbb Q$. The Euler characteristic identity in
the semisimple category of $\mathbb Q[D]$-modules gives the
displayed decompositions. Thus $D$ acts faithfully on
$H_1(N,\mathbb Q)$ and on some finite quotient of the free
abelian group $H_1(N,\mathbb Z)$. But the action of the
nonidentity image of $z$ on every such quotient is induced
by conjugation by $z$, hence is trivial, a contradiction.
The center-free pullback argument of A.3 works continuously:
the kernel of $\operatorname{Aut}(\widehat P)\to
\operatorname{Out}(\widehat P)$ is the compact group
$\operatorname{Inn}(\widehat P)\cong\widehat P$. The two
extensions with the same outer action are therefore isomorphic.
The converse follows by restricting conjugation to the
preserved dense lattice. The finite-index topology argument
also proves that the completion map $E_0\to\widehat E_0$
is injective: surface groups are residually finite and a
nontrivial element outside $P$ is already detected in $Q$.
\end{proof}

\begin{statement}{A.12}{Open}{The unrestricted profinite realization assertion}
We have not proved that every finite outer action on a profinite
hyperbolic loop group admits an invariant discrete surface
lattice, even after replacing that lattice by an isomorphic
discretization. Nor does A.11 prove that every profinite
L-group is the completion of a discrete one. The exact missing
step is a profinite Nielsen realization statement furnishing
the factorization in A.11, up to conjugacy by a continuous
loop-preserving outer automorphism.
\end{statement}
The attack above first uses the specified-lattice criterion,
then the center-free classification of extensions. It proves
realization for all actions obtained by completing discrete
finite actions, but the first step does not apply to an
arbitrary finite subgroup of the profinite outer group.
The referenced LM3 B.9--B.11 concern restriction and extension
of specified discretizations; their non-surjectivity warning
prevents replacing this missing step by an assertion that
every discretization extends. The LM4 finite-actions package
is not among the supplied files. The smallest unresolved
form in this return is a finite-order peripheral outer
automorphism of $\widehat F_2$ with specified hyperbolic
peripheral type: we require either an invariant lattice or
a counterexample to its existence. The finite enumeration
A.10 tests discrete monodromies and cannot decide this
profinite condition. This is an unresolved obligation of
Part A(e), not a completed counterpart theorem. The further attack
in A.16--A.18 proves realization throughout the geometric closure
for the once-punctured torus, and for finite solvable subgroups of
the oriented geometric closure in genus at most two. It also gives
an explicit counterexample if the peripheral hypothesis is removed.
Neither result places an arbitrary peripheral finite subgroup in
that closure.

\begin{statement}{A.13}{Known}{The profinite definitions and peripheral recovery}
For a profinite group $E$, let $\widehat{\mathcal L}_0(R)$
be the structures obtained by identifying $R$ with the full
profinite completion of a hyperbolic surface group, retaining
the completed peripheral homomorphisms and orientation line.
The profinite version of scans 188--193 is
\[
 \widehat{\mathcal L}(E)
   =\mathop{\mathrm{colim}}_{R\triangleleft_o E}
       \widehat{\mathcal L}_0(R)^{E/R}.
\]
All groups and maps are profinite and continuous;
$T$ is a free rank-one $\widehat{\mathbb Z}$-module,
and germs are defined on open submodules $nT$.
The transition maps are injective. The corrected profinite
LT definition is A.2 with a closed normal defining subgroup,
without the open-index requirement, and its specified
commensurability class. The multiplier on $T$ may now be
any element of $\widehat{\mathbb Z}^{\times}$; it is not
restricted by definition to $\{1,-1\}$.
\end{statement}
\begin{proof}
The relevant peripheral recovery theorem says that for an
algebraically simple element $x$ in a nonabelian profinite
surface group and a nonzero integer $n$,
\[
 N(\overline{\langle x^n\rangle})
   =\overline{\langle x\rangle}.
\]
This is \citet[Proposition 3.6]{BoggiZalesskii2019};
its Definition 1.3 includes peripheral simple curves.
Thus a peripheral subgroup is recovered from its intersection
with an open subgroup. Multiplication by a positive integer
is injective on $\widehat{\mathbb Z}$, so its parametrization
is recovered too. The completed orientation line is recovered
by the same degree calculation prime by prime. Passing to
normal open cores proves the displayed filtered definition
has the same germ description. No discrete dense lattice
is part of these data; adding one is the extra assumption
in A.11.
\end{proof}

\begin{statement}{A.14}{Proved}{The centralizer and marking conclusions survive profinitely}
Let $1\to\widehat P\to E\to Q\to1$ be a profinite
hyperbolic L extension and $\rho:Q\to\operatorname{Out}(\widehat P)$
its outer action. Then
\[
 C_E(\widehat P)\cong\ker\rho,
 \qquad
 \ker\rho=1\ \Longleftrightarrow\ C_E(J)=1
       \text{ for every open subgroup }J\leq E.
\]
If a profinite realization is fixed, its continuous marked
isomorphisms form an $\operatorname{Aut}_{L,\mathrm{cts}}(E)$-torsor;
modulo inner automorphisms they form an
$\operatorname{Out}_{L,\mathrm{cts}}(E)$-torsor.
These conclusions do not assert the existence of an
invariant discrete lattice or identify the full continuous
outer group with a discrete orbifold mapping class group.
\end{statement}
\begin{proof}
The homology argument in A.11 also shows
$C_{\widehat P}(U)=1$ for every open $U\leq\widehat P$:
if $z\ne1$ centralizes $U$, choose an open normal
$N\subseteq U$ not containing $z$; the faithful action of
$\widehat P/N$ on the homology of its surface cover gives
the same contradiction. If $\phi\in\operatorname{Aut}(\widehat P)$
fixes an open normal $U$ pointwise, then
$x^{-1}\phi(x)$ centralizes $U$ for every $x$, because
$\phi(xux^{-1})=xux^{-1}$. Thus $\phi=1$.

It follows that an element of $E$ centralizing an open
subgroup of $\widehat P$ centralizes $\widehat P$.
The center-free pullback extension from A.11 identifies
its centralizer with $\ker\rho$. Intersecting any open
$J\leq E$ with $\widehat P$ proves the equivalence,
exactly as in A.4. The torsor assertions are the
composition argument in A.6, using continuous
isomorphisms throughout.
\end{proof}

\begin{statement}{A.15}{Refuted}{The discrete modular automorphism conclusion does not extend literally}
For $E=\widehat{C_2*C_3}$ with its completed cusp structure,
the multiplier map
\[
 \operatorname{Out}_{L,\mathrm{cts}}(E)
       \longrightarrow\widehat{\mathbb Z}^{\times}
\]
is surjective. In particular this outer group is uncountable,
whereas its discrete counterpart in A.8 is $C_2$.
Thus the assertion that all L-automorphisms are the discrete
geometric ones fails profinitely. This does not decide the
finite-subgroup realization problem A.12.
\end{statement}
\begin{proof}
Use the rational stack
$\mathcal X=[(\mathbb P^1-\{0,1,\infty\})/S_3]$, with
the anharmonic $S_3$-action defined over $\mathbb Q$.
Over $\mathbb C$ its signature is $(0;2,3;1)$, so
Riemann existence identifies its geometric \'etale
fundamental group with $E$. The arithmetic fundamental
group sequence gives a continuous outer $G_{\mathbb Q}$
action on $E$. It preserves the cusp germ and acts on
its Tate inertia module $\widehat{\mathbb Z}(1)$ through
the cyclotomic character. This can be checked directly
in a local parameter by the compatible covers
$z\mapsto z^n$: Galois acts on their deck groups by its
action on $\mu_n$. Passing from the surface cusp to the
stack cusp changes only a finite-index subgroup, so
the same multiplier is obtained on germs.

The cyclotomic character is onto
$\widehat{\mathbb Z}^{\times}$, since
$\operatorname{Gal}(\mathbb Q(\mu_n)/\mathbb Q)
 =(\mathbb Z/n)^{\times}$ for every $n$ and restriction
to the union of these fields is surjective. Inner
automorphisms have multiplier one, so these multipliers
are visible in the outer group. This proves the claim.
The arithmetic action used here is exactly the type of
extension anticipated in \S52.2, folio 745, scan 195;
the finite-cover comparison is C.1.
\end{proof}

\subsection{A further attack on unrestricted profinite realization}

The following results address precisely the profinite extension of
the finite-action question in \S52.2, folio 744, scan 193. They
separate a positive theorem about geometric closures from an actual
counterexample without peripheral data; the latter is not presented
as a counterexample to the L-group assertion.

\begin{statement}{A.16}{Refuted}{Finite actions on a bare free profinite group need not be discrete}
Let $\lambda\in\widehat{\mathbb Z}^{\times}$ have component $-1$
in $\mathbb Z_2$ and component $1$ in $\mathbb Z_\ell$ for every
odd prime $\ell$. The assignments
\[
 \alpha(x)=x^\lambda,\qquad \alpha(y)=y
\]
define an involution of $\widehat F_2$. Its semidirect product
$E=\widehat F_2\rtimes_\alpha C_2$ is not the profinite completion
of any finitely generated discrete group. Every induced action on
the free pro-$\ell$ quotient is nonetheless the completion of a
discrete surface action. This involution does not preserve either
standard hyperbolic peripheral structure on the marked rank-two
group; it therefore does not refute A.12.
\end{statement}
\begin{proof}
Profinite exponentiation on the procyclic subgroup generated by $x$
and the universal property of $\widehat F_2$ define $\alpha$.
Since $\lambda^2=1$, substitution twice fixes both generators;
thus $\alpha^2=1$. Its nonidentity action on the abelianization
shows that it has order two, also in the outer group. Abelianizing
the semidirect-product presentation gives
\[
 E^{\mathrm{ab}}\cong
 \widehat{\mathbb Z}\oplus
 \widehat{\mathbb Z}/(\lambda-1)\widehat{\mathbb Z}\oplus C_2,
 \qquad
 \widehat{\mathbb Z}/(\lambda-1)\widehat{\mathbb Z}
 \cong C_2\times\prod_{\ell\ne2}\mathbb Z_\ell.
\]
Consequently its maximal abelian pro-$2$ quotient has free
$\mathbb Z_2$-rank one, whereas every odd-primary quotient has
free rank two. If $E=\widehat G$ for a finitely generated discrete
group, its abelianization would be $\widehat{G^{\mathrm{ab}}}$:
both classify exactly the finite abelian quotients of $G$.
The finitely generated abelian group $G^{\mathrm{ab}}$ has one
integer rank, independent of the prime, a contradiction.
For each odd $\ell$, $\alpha$ is the identity on the free
pro-$\ell$ quotient; for $\ell=2$ it is $x\mapsto x^{-1}$,
$y\mapsto y$, the completion of a reflection of the punctured
torus. The odd-primary actions allow the ineffective $C_2$.
These separate realizations do not supply one full profinite
realization.

Here are finite certificates for the peripheral exclusions.
For the tripod, the three oriented peripheral homology vectors
are $x,y,-x-y$. A peripheral automorphism acts on their rank-two
span as $u\rho(\sigma)$ for a common unit $u$ and the standard
representation of $\sigma\in S_3$. An involution has
$\sigma=1$ or a transposition. In the first case its matrix is
scalar; in the second its trace is zero. The matrix
$\operatorname{diag}(\lambda,1)$ is neither scalar nor of trace
zero. This argument remains valid after any profinite change of
marking, since being scalar and having trace zero are invariant
under conjugation.

For the standard punctured torus put $c=[x,y]$ with commutator
$xyx^{-1}y^{-1}$. If $\alpha(c)$ were conjugate to $c^u$, the
class-two Heisenberg quotients modulo every positive integer
would give $u=\det(\alpha^{\mathrm{ab}})=\lambda$.
Take the finite quotient
\[
 H=(\mathbb Z/4)^2\rtimes\langle t\rangle,
 \quad |t|=3,\quad
 M=\begin{pmatrix}0&-1\\1&-1\end{pmatrix},
 \quad x\mapsto t,\quad y\mapsto(1,0).
\]
The images generate $H$, which has order $48$. The image of $c$
is $(3,1)$ and its conjugacy class is
$\{(3,1),(3,2),(2,1)\}$. Since $\lambda\equiv1\pmod3$,
$\alpha(c)$ has the same image $(3,1)$. Since
$\lambda\equiv-1\pmod4$, $c^\lambda$ has image
$(1,3)$, outside that conjugacy class. Thus the standard torus
peripheral condition fails. This last argument concerns the
specified commutator structure, not every possible new torus
structure on the abstract profinite group.

The program \texttt{code/a\_profinite\_obstructions.py}
constructs this whole finite group, checks all $110592$
associativity triples and all $48$ conjugators, and verifies
the displayed images. It also records the shorter integral
homology obstruction: $\lambda\equiv7\pmod{12}$, so
$\operatorname{tr}(\alpha^{\mathrm{ab}})\equiv8\pmod{12}$;
an integral rank-two involution has trace $-2,0,$ or $2$, whose
residues are $10,0,2$. Thus no change of profinite basis makes
this involution integral on homology.
\end{proof}

\begin{statement}{A.17}{Proved}{Realization for finite subgroups in geometric closures}
Let $\check\Gamma_{1,1}^{\pm}$ be the closure of
$\operatorname{Out}_L(F_2)=\operatorname{GL}_2(\mathbb Z)$
in $\operatorname{Out}_L(\widehat F_2)$ for the
once-punctured-torus structure. Every finite subgroup of this
closure is conjugate inside the closure to a discrete finite
subgroup, and hence every finite outer action with image in
this closure is realized by a finite surface action after a
profinite change of marking. Its effective image is isomorphic
to one of
\[
 1,\ C_2,\ C_3,\ C_4,\ C_6,\ C_2\times C_2,\ S_3,\ D_8,\ D_{12},
\]
where the dihedral subscripts are the group orders. If the
orientation is preserved, only the cyclic groups of orders
$1,2,3,4,6$ occur.

More generally, for every hyperbolic $S_{g,n}$ with $g\le2$,
every finite solvable subgroup of the closure of its
orientation-preserving mapping class group is conjugate
inside that closure to a discrete finite subgroup. The same
realization conclusion therefore holds for those outer
actions. Punctures may be permuted. These hypotheses concern
the geometric closure, not the full peripheral outer group.
\end{statement}
\begin{proof}
The Main Theorem of \citet[p.~2 of the primary arXiv version]{BuxErshovRapinchuk2011}
asserts exactly that
$\widehat{\operatorname{Out}(F_2)}\to
\operatorname{Out}(\widehat F_2)$ is injective; this is the
outer form of Asada's congruence theorem. Compactness identifies
its image with the geometric closure.
The standard modular amalgam, also recorded explicitly by
\citet[p.~3]{Zimmermann2021}, is
\[
 \operatorname{SL}_2(\mathbb Z)=C_4*_{C_2}C_6,
 \qquad
 \operatorname{GL}_2(\mathbb Z)=D_8*_{C_2\times C_2}D_{12}.
\]
For explicit vertex groups take
\[
 S=\begin{pmatrix}0&-1\\1&0\end{pmatrix},\quad
 U=\begin{pmatrix}0&-1\\1&1\end{pmatrix},\quad
 R=\begin{pmatrix}0&1\\1&0\end{pmatrix}.
\]
Then $S^2=U^3=-I$, $S^4=U^6=R^2=I$, and $R$
inverts both $S$ and $U$. Thus the second amalgam follows
from the first by adjoining this reflection. Completion
gives the profinite amalgam with the same finite vertex and
edge groups: this follows directly from their universal
mapping properties for maps to finite groups. The vertex
maps are injective because the discrete matrix group is
residually finite, as detected by reduction modulo integers;
the edge inclusions are proper.

The exact result now used is
\citet[Theorem 3.10, p.~421]{ZalesskiiMelnikov1989}:
every finite subgroup of the fundamental profinite group of
a finite graph of profinite groups satisfying the injectivity
and nondegeneracy assumptions of its \S3.6 is conjugate into
a vertex group. The preceding paragraph checks both
assumptions for our amalgams. Hence a finite subgroup of
$\check\Gamma_{1,1}^{\pm}$ is conjugate into $D_8$ or
$D_{12}$, and an oriented one into $C_4$ or $C_6$.
This is a theorem about the entire closure, not merely
finite subgroups already lying in the discrete group.
The exact code exhausts the $10$ subgroups of the first
dihedral vertex group and the $16$ of the second, producing
precisely the nine abstract types displayed above. These
counts are subgroup counts, not conjugacy-class counts.

For the general low-genus assertion, the known congruence
property for $g\le2$ identifies the geometric closure with
the full profinite completion of the mapping class group
(the scope and references are stated in
\citealp[\S1.1]{Boggi2026Symmetric}). The relevant newer
finite-subgroup theorem is
\citet[Theorem C(i), equivalently Theorem 5.1(i) with empty boundary]{Boggi2026Symmetric}:
the inclusion of that discrete mapping class group into its
completion gives a bijection on conjugacy classes of
\emph{finite solvable} subgroups. The definition in its
\S4 uses orientation-preserving diffeomorphisms; ``full''
allows boundary or puncture permutations and does not add
orientation reversal. Thus the theorem applies with exactly
the convention stated here. A finite subgroup of the closure
lifts uniquely to a finite subgroup of the completion and
is conjugate to a discrete one by this theorem.

Finally, for an outer action $\rho:Q\to\check\Gamma$, choose
the conjugating element just obtained. The conjugated
homomorphism factors through the discrete peripheral outer
group. A.3 realizes it and A.11 identifies the completed
extension; transporting the marking back realizes the
original profinite extension. Any kernel of $\rho$ remains
the ineffective kernel, so no assertion bounds $|Q|$ when
ineffective actions are permitted.
\end{proof}

\begin{statement}{A.18}{Proved}{Necessary orientation tests and the remaining theorem obstruction}
A finite peripheral action which is geometrically realizable
must have orientation multiplier in the diagonally embedded
$\{1,-1\}\subset\widehat{\mathbb Z}^{\times}$. On homology
its characteristic polynomials must come from integral
polynomials independently of the prime. In tripod type,
every finite subgroup fixing all three peripheral classes
must then be trivial or $C_2$; in the second case its
nonidentity element has multiplier $-1$.

Consequently A.12 in tripod type would imply that every
finite subgroup of the profinite Grothendieck--Teichm\"uller
group, viewed in the pure peripheral outer group, is trivial
or $C_2$, and its multiplier-one subgroup has no nontrivial
torsion. The checked finite-subgroup and rigidity theorems
do not establish this conclusion for the full peripheral
outer group.
\end{statement}
\begin{proof}
Any discrete surface homeomorphism acts on its integral
orientation line by $1$ or $-1$, and on its integral first
homology by an integral matrix. Completion preserves both
facts. Changing the profinite marking conjugates the
homology matrices and leaves the orientation multiplier
unchanged, since that multiplier takes values in an abelian
group. Thus these are necessary conditions even when the
discretization may change. They are not asserted sufficient.
The full discrete mapping class group of the tripod is
$S_3\times C_2$, with pure subgroup $C_2$ generated by the
reflection fixing $0,1,\infty$; this follows by extending
over the three punctures, isotoping a three-point-fixing
orientation-preserving sphere map to the identity, and
using complex conjugation for the other orientation.
Permutation of peripheral classes commutes with passage
to a new peripheral marking. Hence a pure finite subgroup
which is realized lies in that pure $C_2$.

For precision, the usual $\widehat{GT}$ action has the
form $x\mapsto x^u$, $y\mapsto f^{-1}y^u f$, where
$f\in[\widehat F_2,\widehat F_2]$; its defining
hexagon relations ensure the same powered-conjugacy
condition for $(xy)^{-1}$. Its image is therefore pure
peripheral. This action is also faithful in the outer
group. Indeed, if it is conjugation by $g$, abelianization
gives $u=1$. Then $g$ centralizes $x$, so $g=x^a$ by
the peripheral centralizer theorem in A.13. The equality
on $y$ gives $fx^a=y^b$ for some $b$. Abelianizing and
using $f^{\mathrm{ab}}=0$ forces $a=b=0$, then $f=1$.
Applying the preceding pure-subgroup conclusion proves
the stated implication of A.12. This implication concerns
conjugacy in the full peripheral outer group; it does not
claim a conjugator inside $\widehat{GT}$.

Here are the exact boundaries of the literature routes
checked in this attack. \citet[Theorem A and Theorem 6.2]{BoggiZalesskii2025}
treat finite $p$-subgroups in completions of specified good
discrete groups, including hyperelliptic mapping class
groups; their Theorem C adds hypotheses on normalizers
to handle finite solvable subgroups. They do not start
with an arbitrary subgroup of
$\operatorname{Out}_L(\widehat P)$.
\citet[Theorems A, 4.1 and 5.1]{Boggi2026Symmetric}
starts its centralizer argument with a finite action
already on a surface, and its conjugacy conclusion is
inside the mapping class completion, as used in A.17.
\citet[Theorem 1.1]{HerfortZalesskii2013} proves that a
finitely generated virtually free \emph{pro-$p$} group
is the fundamental pro-$p$ group of a finite graph of
finite $p$-groups. Its conclusion cannot be transferred
to the full profinite category: A.16 is an explicit
counterexample to that transfer. Finally, the 2026
\citet[Theorem 1.1]{LochakNakamuraSchneps2026} constructs
a $\widehat{GT}$ action on completed mapping class
groups; it does not classify torsion inside
$\widehat{GT}$ or construct invariant discrete lattices.

For historical calibration, the question about torsion
in $\widehat{GT}$ beyond conjugates of complex conjugation
was explicitly posed by
\citet[\S1.4(4), p.~190]{Schneps1997Survey}. We use that
passage only to identify the mathematical obstruction,
not as evidence that every version remains open in 2026,
and not as a substitute for the full peripheral hypothesis.
The missing step after A.17 remains to put an arbitrary
finite peripheral subgroup, possibly after peripheral
conjugacy, into a geometric closure, or to exhibit a
genuinely peripheral counterexample. Neither membership
nor a finite-order lift of a mixed cyclotomic multiplier
has been obtained here.
\end{proof}
